% =====================================================================
% Plain-English explainer for the research paper.
%
% Companion to jcst_paper.pdf. Written for the authors to read before
% presenting or defending the work, and for anyone marking it who wants
% the argument without the apparatus.
%
% Compile:  tectonic -X compile explainer.tex
% =====================================================================

\documentclass[11pt]{article}

\usepackage[a4paper,top=0.95in,bottom=0.95in,left=1.05in,right=1.05in]{geometry}
\usepackage{newtxtext}
\usepackage{newtxmath}
\usepackage{booktabs}
\usepackage{array}
\usepackage{enumitem}
\usepackage{xcolor}
\usepackage{titlesec}
\usepackage[hidelinks]{hyperref}
\usepackage{parskip}
\usepackage{fancyhdr}
\usepackage{tcolorbox}

\definecolor{accent}{HTML}{1D4ED8}
\definecolor{warn}{HTML}{B45309}
\definecolor{rule}{HTML}{CBD5E1}

\titleformat{\section}{\large\bfseries\color{accent}}{\thesection.}{0.5em}{}
\titleformat{\subsection}{\normalsize\bfseries}{\thesubsection}{0.5em}{}
\titlespacing{\section}{0pt}{14pt}{5pt}
\titlespacing{\subsection}{0pt}{10pt}{3pt}

\setlist[itemize]{leftmargin=1.3em,itemsep=2pt,topsep=3pt}
\setlist[enumerate]{leftmargin=1.5em,itemsep=3pt,topsep=3pt}

\pagestyle{fancy}
\fancyhf{}
\fancyhead[L]{\footnotesize\color{accent}The paper, in plain English}
\fancyhead[R]{\footnotesize\thepage}
\renewcommand{\headrulewidth}{0.4pt}
\renewcommand{\headrule}{\color{rule}\hrule height 0.4pt}

\newtcolorbox{callout}[1][]{
  colback=warn!5, colframe=warn!60, boxrule=0.6pt, arc=2pt,
  left=7pt, right=7pt, top=5pt, bottom=5pt, #1}

\begin{document}

\begin{center}
{\LARGE\bfseries The paper, in plain English\par}
\vspace{5pt}
{\large Deterministic Guardrails and Variance-Aware Evaluation\\
for Multi-Agent Content Personalization\par}
\vspace{8pt}
{\normalsize Prachi Sharma \quad Danish Kaul \quad Dishant Sherawat\par}
{\footnotesize School of Engineering and Technology, Sushant University,
Gurugram, Haryana\par}
\vspace{6pt}
{\footnotesize\itshape A companion to the research paper. No jargon that
is not explained.\par}
\end{center}

\vspace{6pt}
\hrule height 0.4pt
\vspace{4pt}

% ---------------------------------------------------------------------
\section{Is the paper about our project?}

\textbf{Yes, entirely.} It is not a literature survey or a side essay. Every
number in it was produced by running our own code, and every section describes
something we actually built.

\begin{center}
\small
\begin{tabular}{@{}llr@{}}
\toprule
Paper section & Code that implements it & Size\\
\midrule
3.1 Eight-stage pipeline   & \texttt{app/agents/graph.py} & 461 lines\\
3.2 Three-level profiles   & \texttt{app/services/personalization.py} & 466 lines\\
3.3 Scope guard            & \texttt{app/agents/scope.py} & 138 lines\\
3.4 Grounding guard        & \texttt{app/agents/grounding.py} & 178 lines\\
4.1 The test dataset       & \texttt{marketing\_eval/cases.toml} & 30 cases\\
4.2 Automatic checks       & \texttt{marketing\_eval/harness.py} & 110 assertions\\
4.3 Rubric and judge       & \texttt{rubric.md}, \texttt{judge.py} & 32 dimensions\\
5.3 Model comparison       & \texttt{compare\_models.py} & 7 runs\\
5.6 The empty-reply defect & \texttt{llm/openai\_compatible.py} & in \texttt{generate()}\\
\bottomrule
\end{tabular}
\end{center}

The system is 76 Python modules, roughly 12{,}600 lines of application code, and
10 verification scripts. One of those scripts, \texttt{check\_claims.py},
re-reads the stored run reports and confirms that all 76 numbers in the paper
still match them. So if anyone asks where a figure such as 4.28 comes from, the
answer is a file on disk, not somebody's memory.

% ---------------------------------------------------------------------
\section{What problem were we solving?}

Ask a language model to write marketing copy and two things go wrong.

\textbf{It does not know your brand.} It writes in a generic voice, picks the
wrong conventions for the platform, and ignores the fact that you never use
emojis.

\textbf{It invents numbers.} Ask for a post about cloud costs and it will
cheerfully write ``cut your bill by 40\%'' --- a figure nobody gave it. In
marketing copy that is not a harmless slip. It is a false claim published under
your brand name.

\begin{callout}
We tried the obvious fix for the second problem first: we told the model, in its
instructions, not to invent figures. \textbf{It kept inventing figures.} That
failure shaped the entire project.
\end{callout}

% ---------------------------------------------------------------------
\section{What we built}

Five narrow agents rather than one general chatbot: a caption writer, a
search-optimization agent, a campaign planner, and two that produce images (a
social post image and an advertisement creative). Each has its own instructions
and its own settings.

Three parts are worth understanding properly.

\subsection{Three-level profiles}

Your details live at the level where they are actually true.

\begin{itemize}
\item \textbf{Global} --- who you are, how you like to be spoken to.
\item \textbf{Domain} (marketing) --- brand name, industry, audience, tone,
      banned words.
\item \textbf{Agent} --- for the caption writer: which platforms, how long,
      hashtag policy.
\end{itemize}

State your brand tone once and all five agents receive it. State ``at most three
hashtags'' and only the caption writer receives it, because it is meaningless to
the others.

\subsection{The scope guard}

If you are in the caption writer and you ask for an advertisement image, the
system \textbf{hands you off} to the right agent instead of doing a poor job.

The important part: that decision is made in \textbf{plain code, with no model
call.} It is instant and free. It compares how much evidence your request gives
for the current agent's job against every other agent's job, and redirects only
when the mismatch is obvious --- four separate conditions must all hold, because
wrongly redirecting somebody is worse than answering slightly off-topic.

\begin{callout}
\textbf{A bug worth mentioning to an examiner.} The first version scored each
agent using its whole vocabulary \emph{including its own name}. So ``design an
ad creative'' sent to the caption writer scored highly \emph{for the caption
writer} --- because the phrase contains words from its own name --- and was
never redirected. The fix was to score only \emph{task} words and ignore name
matches.
\end{callout}

\subsection{The grounding guard}

After the model writes something, we scan the text for claim-shaped numbers ---
percentages, money, ``10x'', ``up to'' --- and check each one against the
numbers you actually supplied in your request, your attachment or your profile.
Anything we cannot account for is \textbf{flagged}.

Two design decisions matter here.

\textbf{It flags; it never rewrites.} Deleting a number mid-sentence leaves you
with ``reduce costs by up to'', which is worse than the original. And only
\emph{you} know whether the claim is true, so the system's job is to tell you,
not to decide for you.

\textbf{It checks provenance, not truth.} We are not fact-checking the world. We
are asking a much cheaper and more useful question: did this number come from
the user, or did the model make it up?

% ---------------------------------------------------------------------
\section{How we measured it, and why that is the real contribution}

This is the part a reviewer cares about most.

\subsection{Two kinds of question}

Some questions a computer can answer exactly:

\begin{itemize}
\item Did it hand off to the right agent?
\item Did the budget split add up to the 400 the user gave?
\item Did it stay under three hashtags?
\item Is the generated image exactly 1200$\times$627 pixels?
\end{itemize}

We wrote \textbf{110 such checks} across 30 test cases. They run in seconds,
cost nothing, and give the same answer every time.

Other questions a computer cannot answer: is this a good opening line? Does it
sound like our brand? Would you actually publish it? For those we used an
\textbf{LLM judge} --- a second, independent AI model that reads the output
against a written scoring guide (32 dimensions, scored 1--5) and returns a score
with a reason.

\subsection{Then we checked whether the judge could be trusted}

Most people stop at ``the judge gave us 4.23 out of 5'' and report it. We did
not, and this is where the interesting results came from.

\textbf{We asked the judge the same question three times} --- same outputs, same
prompt, temperature set to zero, so it should be deterministic. It disagreed
with itself, by a mean spread of \textbf{0.51 to 0.54 points on a five-point
scale}.

\textbf{Then we re-ran the whole system without changing anything} and judged
again. One agent's score moved by \textbf{0.37} with zero changes to the code.

\begin{callout}
That single number reframed everything. It means any quality difference smaller
than about half a point is \textbf{noise, not a result.}
\end{callout}

% ---------------------------------------------------------------------
\section{What we found}

\subsection{Finding 1: the guards worked}

Holding the model fixed, adding the guards took the automatic test suite from
\textbf{102/105 to 105/108} passing checks. Three failures remained, and no
amount of prompt rewriting fixed them.

We then swapped in a bigger model and changed nothing else. All three passed and
the suite reached \textbf{110/110}. So those three failures were the
\emph{model's} limitation, not our design's. You can only make that claim if you
change one thing at a time, which is what a controlled experiment is for.

\subsection{Finding 2: our scoring scale had silently broken}

The rubric defines what a 1, a 3 and a 5 look like for each dimension. We
assumed the judge would use 2 and 4 for work that fell in between.

It never did. Across \textbf{174 scores, every single one was a 1, a 3 or a 5.}
Not one 2, not one 4. Our five-point scale had quietly become a three-point
scale, and nothing in the system would ever have told us --- we saw it only
because we plotted the distribution.

Naming the middle values in the judge's instructions brought them back and cut
the worst disagreement from 4 points to 2.

\subsection{Finding 3: ``best model for each agent'' sounds clever and is not}

We had five specialised agents, so the natural next step was to give each one
the model best suited to it. A first screening pass appeared to confirm it: two
agents looked 0.45 and 0.33 better under different models.

So we built it properly and tested it against a control.

\begin{center}
\small
\begin{tabular}{@{}lcc@{}}
\toprule
 & One model for all & Mixed fleet\\
\midrule
Overall score & \textbf{4.28} & 4.18\\
Speed & \textbf{54 s} & 201 s\\
The two ``improved'' agents & 4.56 / 3.80 & 4.33 / \textbf{3.87}\\
\bottomrule
\end{tabular}
\end{center}

The mixed fleet was \textbf{worse overall and 3.7 times slower}, and one of the
two ``improved'' agents actually got \emph{worse}. A proper statistical test
(Wilcoxon signed-rank, 15 changed pairs) gave $p = 0.42$, nowhere near
significant.

Why did screening mislead us? Because 0.33 and 0.45 sit \emph{inside} the noise
floor we had measured. They were never real.

\begin{callout}
We report this as a failure, in full, in the paper. That is deliberate. A
negative result that saves somebody else three days of work is worth more than a
positive result we had to squint at.
\end{callout}

\subsection{Finding 4: a bug that looked like somebody else's fault}

While comparing models, four of them failed on between 4 and 19 test cases each,
reporting what looked like the provider being down.

It was not. Some newer models ``think'' before answering, and that thinking
consumes the token budget. When it runs out mid-thought, the interface returns
\textbf{success, with an empty message}. Our code saw an empty message and
reported a provider outage. One model returned 700 tokens of pure thinking and
\emph{zero} characters of answer, on one request in six.

We now detect that exact signature and retry with more room. Five of six models
went to \textbf{zero errors}. The sixth is genuinely unusable for this task: it
produced 12{,}000 to 16{,}000 characters of reasoning and never answered,
needing 41 retries in a single run.

This mattered beyond robustness. Before the fix, each model failed on a
\emph{different} set of cases, so each model's score described a different
subset of the test. Any ranking would have been an accident of which cases
happened to survive.

% ---------------------------------------------------------------------
\section{The one-sentence version of each takeaway}

\begin{enumerate}
\item \textbf{If you can check something in code, check it in code.} Telling a
      model ``do not invent numbers'' does not work. A 178-line checker does.
\item \textbf{Flag, do not fix.} When the system cannot know whether something
      is wrong, tell the user instead of silently editing.
\item \textbf{Publish your noise floor, not just your score.} We measured ours
      at half a point, and it invalidated two conclusions we had already drawn.
\item \textbf{Specialised agents do not need specialised models.} We tested it;
      the difference was smaller than the noise and cost 3.7 times the latency.
\item \textbf{Test your test code.} Both of our worst bugs were in the
      \emph{measuring} equipment, and both produced confident, wrong numbers.
\end{enumerate}

% ---------------------------------------------------------------------
\section{Honest limitations}

Expect to be asked about these. They are stated plainly in the paper's Threats
to Validity section, which is the correct place for them.

\begin{itemize}
\item \textbf{No human scoring pass.} Our judge was never checked against real
      human raters. This is the biggest weakness, and we say so.
\item \textbf{One brand, one language, one industry.} All results use a single
      synthetic brand in English.
\item \textbf{We wrote the 30 test cases ourselves}, to probe failures we
      expected. They are not sampled from real traffic.
\item \textbf{Small sample for the statistics.} Fifteen changed pairs. A larger
      dataset might find an effect we could not.
\item \textbf{The noise floor of 0.5 is ours}, for this task and this judge. It
      is not a universal constant.
\end{itemize}

% ---------------------------------------------------------------------
\section{Likely questions, and how to answer them}

\textbf{Why not just fine-tune a model on the brand?}\\
Too expensive per user, and it cannot be updated the moment a brand changes its
tone. Composing the profile at request time is instant and works per user.

\vspace{5pt}
\noindent\textbf{Is using an AI to judge an AI not circular?}\\
Yes, which is why we measured how much it disagrees with itself, used a
\emph{different} model as judge than the one being judged, and made the code
flag loudly whenever those two are the same. Section 2.4 of the paper cites the
published work on self-preference bias.

\vspace{5pt}
\noindent\textbf{Why is a negative result in the paper?}\\
Because it is true, and because the reasoning behind it --- effect size against
measured noise --- is the paper's main methodological point. Removing it would
leave out the most useful finding.

\vspace{5pt}
\noindent\textbf{What is actually new here?}\\
Three things: scope enforcement as a deterministic, zero-cost step rather than a
prompt instruction; provenance-checking of claim-shaped numbers rather than
fact-checking; and reporting judge variance as a required part of judge-based
evaluation.

\vspace{5pt}
\noindent\textbf{Could you have used a bigger model and skipped all this?}\\
A bigger model fixed three test failures, and we show that. It did \emph{not}
stop figure fabrication, and it cannot perform the scope handoff for free. The
guards and the model solve different problems.

\end{document}
